%% %%%%%%%%%%%%%%%%%%%%%%%%%%%%%%%%%%%%%%%%%%%%%%%%%
%% Template for working papers
%% Tommaso Di Francesco
%% University of Amsterdam
%% %%%%%%%%%%%%%%%%%%%%%%%%%%%%%%%%%%%%%%%%%%%%%%%%%
%% Version 1.0 // February 2024
%% 
\documentclass[11pt]{article}


\usepackage[utf8]{inputenc}
\usepackage[T1]{fontenc}
\usepackage{geometry}
\usepackage{graphicx}
\usepackage{amsmath}
\usepackage{amssymb}          % For \mathbb
\usepackage{amsthm}           % For theorem environments
\usepackage[backend=biber,style=authoryear,natbib=true,maxcitenames=2,maxbibnames=99]{biblatex}
\addbibresource{references.bib}
\usepackage{setspace}         % For \onehalfspacing
\usepackage{titling}          % For \droptitle
\usepackage{xcolor}           % For \textcolor
\usepackage{hyperref}



\geometry{
    a4paper,
    left=2.5cm,
    right=2.5cm,
    top=2.5cm,
    bottom=2.5cm
}

% Define theorem environments
\newtheorem{assumption}{Assumption}
\newtheorem{proposition}{Proposition}
\newtheorem{definition}{Definition}
\newtheorem{remark}{Remark}

\onehalfspacing
%% ===============================================

\setlength{\droptitle}{-5em} %% Don't touch

% %%%%%%%%%%%%%%%%%%%%%%%%%%%%%%%%%%%%%%%%%%%%%%%%%%%%%%%%%%
% SET THE TITLE
% %%%%%%%%%%%%%%%%%%%%%%%%%%%%%%%%%%%%%%%%%%%%%%%%%%%%%%%%%%
% TITLE:
\title{Lasso Trading}
\author{}
% DATE:
\date{\today}

% %%%%%%%%%%%%%%%%%%%%%%%%%%%%%%%%%%%%%%%%%%%%%%%%%%%%%%%%%%
% %%%%%%%%%%%%%%%%%%%%%%%%%%%%%%%%%%%%%%%%%%%%%%%%%%%%%%%%%%
\begin{document}
% %%%%%%%%%%%%%%%%%%%%%%%%%%%%%%%%%%%%%%%%%%%%%%%%%%%%%%%%%%
% %%%%%%%%%%%%%%%%%%%%%%%%%%%%%%%%%%%%%%%%%%%%%%%%%%%%%%%%%%

{\setstretch{.8}
\maketitle
\centering
}

\section{Introduction}


The volume and granularity of data available worldwide are growing at an unprecedented pace.
In 2024 alone, the world generated approximately 402 million terabytes of data per day, up from 43 million just a decade earlier \citep{DataSphere2025}.
Whereas investors once had access to only a limited set of financial time series, they can now draw on a vast number of signals, including balance-sheet information, textual data, satellite imagery, and web traffic.
The way in which investors process this data is fundamentally different from most asset-pricing models where agents are endowed with a theoretical model featuring a prespecified set of signals that are believed to contain information for future returns.
Agents then use econometric methods to estimate the models parameters and form return expectations.
Recent survey evidence shows that this is ceasing to be a good approximation of real investor behavior.
Instead of relying on prespecified theoretical models, a majority of asset managers now use machine learning and artificial intelligence to process large amounts of data \citep{EY2025_GenAI_WealthAssetMgmt, GrantThornton2025_AIAssetMgmt}. 
In other words, the way investors learn about market behavior has shifted from theory-driven to data-driven.

This shift is not well captured by existing learning-based asset-pricing models.
The common feature of this model class is that agents eihter know the true model which returns follow or that they hold a misspcified model that is either under- or overparameterized.
In both cases, the set of predictors is small and fixed.
We argue that this framework is ill-suited to study modern learning behavior in financial markets.
Instead, we propose 

\subsection{Relation to the Literature}





While classical learning-based asset-pricing models discipline beliefs through low-dimensional perceived laws of motion (PLMs) and recursive least squares updating, modern systematic investors routinely confront high-dimensional predictor spaces with limited theoretical guidance.
In such environments, sparsity-inducing regularization methods---most prominently the least absolute shrinkage and selection operator (LASSO)---are widely used to select predictors and generate forecasts in finite samples.

This paper studies the equilibrium implications of such behavior.
Specifically, we ask: \emph{What are the asset-pricing consequences of investors who actively search for return predictability using high-dimensional, regularized learning rules, even when returns are unpredictable under rational expectations?}

We develop an equilibrium mechanism we term \emph{Lasso Trading}.
In the model, investors do not possess a structural theory linking fundamentals to returns.
Instead, they entertain the possibility that (log) returns are conditionally predictable from a large set of candidate predictors and estimate forecasting rules using LASSO \citep{Tibshirani1996}.
The central force is a feedback loop from learning to prices: the act of searching for predictability affects trading decisions and hence equilibrium returns.
As a result, the economy endogenously generates \emph{sparse and short-lived episodes of return predictability}, despite the absence of predictability under rational expectations.

The mechanism speaks directly to empirical evidence showing that return predictability is episodic, unstable, and concentrated in a small, shifting subset of predictors \citep{Chinco2019,Farmer2023}.
Our contribution is twofold.
First, we provide, to our knowledge, the first analytical characterization of LASSO-based learning in an equilibrium asset-pricing environment, thereby linking the macro-finance learning literature to the empirical tradition of high-dimensional return forecasting.
Second, the model delivers sharp, testable implications: predictability is (i) sparse, (ii) unstable over time, and (iii) more likely when estimation relies on short memory, weak penalization, and high return volatility.
We bring these predictions to the data using a large predictor set that includes text-based topics extracted from financial news \citep{Bybee2021}, and we estimate a structural feedback parameter governing how learning-induced beliefs map into realized returns.

\subsection*{Relation to the Literature}

The paper relates to three strands of research.

\paragraph{Learning in asset pricing.}
A large literature studies bounded rationality by assuming that agents learn unknown parameters of a correctly specified model using least squares methods \citep{Marcet1989}.
Classic results show that learning can generate excess volatility and return predictability \citep{Timmermann1993,Timmermann1996}.
Subsequent work emphasizes that convergence to rational expectations need not obtain when agents have limited memory or employ misspecified forecasting models \citep{Honkapohja2003}.
Related contributions allow agents to switch among multiple forecasting rules or to rely on parsimonious PLMs in multivariate environments \citep{barberis1998model,hong2007simple,Branch2010}.

\paragraph{Internal rationality and belief--price feedback.}
A central challenge in standard learning models is that agents often know the mapping from fundamentals to prices.
Following \cite{Adam2016}, we adopt an internal rationality perspective in which agents optimize conditional on subjective beliefs about returns and learn from realized prices.
This framework generates a feedback loop between beliefs and equilibrium outcomes.
We retain the analytical discipline of a Lucas-tree economy \citep{Lucas1978}, while replacing ad hoc learning rules with a microfounded, empirically prevalent high-dimensional procedure.

\paragraph{Sparse and unstable predictability.}
Empirical evidence documents that predictive relations are time-varying \citep{AngBekaert2007_RFS,GoyalWelch2003_MS,LettauLudvigson2001_JF} and sparse across a large set of candidate signals \citep{kozak2020shrinking,freyberger2020dissecting,Feng2020}.
Machine-learning methods exploit this sparsity and often outperform traditional linear models \citep{gu2020empirical}.
Most closely related, \cite{Chinco2019} show that LASSO uncovers high-frequency predictability that is unstable and concentrated in a small set of predictors.
Our contribution is to endogenize these features as equilibrium outcomes of learning and trading.

\subsection*{Model}

\paragraph{Environment.}
We study a Lucas-tree economy following \cite{Adam2016}.
A representative risky asset pays dividend $D_t$ and satisfies the Euler equation
\begin{equation}
P_t = \delta\,\mathbb{E}_t\!\left[\left(\frac{C_{t+1}}{C_t}\right)^{-\gamma}(P_{t+1}+D_{t+1})\right],
\end{equation}
where preferences exhibit constant relative risk aversion with coefficient $\gamma$ and discount factor $\delta$.
Under rational expectations, consumption and dividends follow stationary log-normal processes and returns are conditionally unpredictable.

\paragraph{Beliefs and perceived law of motion.}
Agents are fully rational with respect to the dividend and consumption processes but are boundedly rational with respect to returns.
They entertain a linear perceived law of motion for log returns,
\begin{equation}
r_{t+1} = X_t'\beta_t + \eta_{t+1},
\end{equation}
where $X_t$ is a potentially high-dimensional vector of predictors.
Agents treat the predictor set as a large menu of candidate signals and estimate $\beta_t$ from the data.

\paragraph{Bounded expected returns.}
To ensure global existence of equilibrium, we bound perceived expected gross returns using a smooth saturation function,
\begin{equation}
\mathbb{E}\!\left(R_{t+1}\mid \mathcal{F}_t\right)
= \phi\,\exp\!\Bigl(w\tanh(X_t'\beta_t / w)\Bigr),
\end{equation}
which implies a well-defined price-dividend ratio and an actual law of motion for returns featuring
\[
g_{\beta}(x) := \log\!\left(1-\kappa\,e^{w\tanh(x\beta/w)}\right).
\]

\paragraph{Equilibrium concept.}
Because the perceived model is misspecified relative to the nonlinear equilibrium mapping, we adopt a moment-based equilibrium concept.
Let $\mathcal{R}$ denote the regression operator used by agents.
A \emph{Cross-Moment Consistent Equilibrium} is a pair $(\mu,\beta)$ such that (i) $\mu$ is invariant under the actual law of motion implied by $\beta$, and (ii) $\beta$ equals the population regression coefficient produced by $\mathcal{R}$ under $\mu$.

\subsection*{Learning and the Main Result}

We study three learning schemes: decreasing-gain recursive least squares, constant-gain weighted least squares, and LASSO.

Under decreasing-gain learning, the unique equilibrium satisfies $\beta=0$ and coincides with rational expectations.
With constant-gain learning, parameter estimates fluctuate persistently around zero but do not converge.
Introducing LASSO regularization induces sparsity by thresholding small coefficients.

In the standardized benchmark, the LASSO estimator satisfies
\begin{equation}
\hat{\beta}_{\text{LASSO}}
=\operatorname{sign}(\hat{\beta}_{\text{OLS}})
\max\{0,|\hat{\beta}_{\text{OLS}}|-\lambda\}.
\end{equation}
Hence, agents trade on a predictor only when recent data generate sufficiently strong estimated predictability.

\paragraph{Main implication.}
Even though the equilibrium value satisfies $\beta=0$, finite-sample learning with short memory and regularization generates temporary, sparse deviations from zero.
The probability that a predictor is selected is increasing in return volatility and decreasing in the penalty parameter.
These deviations constitute endogenous windows of predictability.

\paragraph{Economic mechanism.}
Lasso Trading generates predictability through a belief--price feedback loop.
Agents estimate predictability from recent data, trade on perceived signals, and thereby affect realized returns.
Regularization ensures sparsity, while short memory ensures instability.
Predictability is therefore episodic, fragile, and migrates across predictors over time.

\subsection*{Empirical Strategy}

We implement a two-stage empirical design.
In the first stage, we estimate rolling LASSO forecasting regressions using a high-dimensional predictor set that includes financial variables and news-based text topics.
In the second stage, we estimate the equilibrium feedback parameter $\kappa$ from the nonlinear return mapping implied by the model.

The theory yields sharp predictions regarding sparsity, instability, and the conditions under which learning-driven trading has economically meaningful price effects.

\subsection*{Conclusion}

This paper provides a tractable equilibrium framework in which high-dimensional, sparsity-inducing learning rules endogenously generate sparse and transient return predictability.
The model rationalizes empirical patterns documented in the data and connects modern machine-learning forecasting practices to equilibrium asset pricing.

\newpage

\printbibliography
\end{document}